% MT024 / explanation — canonical editable content.
这里\(X,Y\)是已有训练输入与标签，\(w\)是要学习的权重。虽然原题把数据损失简写成\(L(X,Y)\)，它实际上还依赖\(w\)产生的预测。整个目标同时要求“解释好训练数据”和“不要使用不必要的大权重”。


\textbf{(a) 平方范数约束权重。} \(\|w\|^2=\sum_jw_j^2\)，权重越大代价越大。它使模型在降低训练损失的收益与增加复杂度的代价之间取舍，有助于减少过拟合；并不要求每个权重最终都相同。


\textbf{(b) 用未参与该次拟合的数据选\(\alpha\)。} 它是控制取舍的超参数。补充例子：对\(\alpha=0.01,0.1,1,10\)分别完成训练内CV，比较平均验证表现。选好一个强度后，用全部训练数据再拟合权重。不能直接以训练损失最低来选，否则常会偏向限制最弱的模型；测试集也不参与选择。


如果把\(\alpha\)也放进原训练目标中任其最小化，且\(\alpha\ge0\)，它可能直接取0以去掉惩罚，失去用新数据检查泛化的意义。超参数与权重的学习职责不同，完整CV流程见\StudyRef{Answers}{mt004}{MT004}。


\textbf{(c) L1鼓励精确为0的系数。} 用\(\sum_j|w_j|\)替换平方和后，一些特征可能完全退出预测式；L2通常只是把它们缩小。L1能用于特征选择，但留下多少特征取决于数据与强度，不是固定比例。


特征尺度会影响两种惩罚的比较：同一个量换单位，表示它所需的权重会变。因此通常先在训练数据上确定标准化参数，再用于验证或测试。


\textbf{出处：} 2021A期中 Q11，原卷第4页（10分）


\textbf{答案依据：} 2021A期中 Q11 随卷／配套参考答案；条件与解释经复核。
